\documentclass[10pt,a4paper]{amsart}
\usepackage[margin=19mm]{geometry}
\usepackage{amsmath,amssymb,mathtools}
\usepackage[scr=rsfs,cal=euler]{mathalfa}
\usepackage{xfrac}
\usepackage[unicode,hidelinks,bookmarks=false]{hyperref}
\newcommand{\ie}{i.e.\ }
\newcommand{\cf}{cf.\ }
\newcommand{\resp}{respectively\ }
\newcommand{\Subsection}{\subsection}
\providecommand{\nobreakdash}{\nobreak\hbox{-}}
\setlength{\emergencystretch}{2em}
\multlinegap0pt
\usepackage{amssymb}
\usepackage{mathtools}
\usepackage{bm}
\usepackage[shortlabels]{enumitem}
\newtheorem{proposition}{Proposition}
\newcommand{\e}{\mathrm e}
\title{\textbf{Bargmann--Fock Representation and Global Estimates\\ for the Linearized Hard-Sphere Boltzmann Operator}}
\author{Ilya Karlin}
\date{Source: arXiv:2609.09548v1 (CC BY 4.0)}
\begin{document}
\maketitle

\noindent\emph{Source and licence.} \textbf{Bargmann--Fock Representation and Global Estimates\\ for the Linearized Hard-Sphere Boltzmann Operator}. Version arXiv:2609.09548v1 (CC BY 4.0), distributed by arXiv under the Creative Commons Attribution 4.0 International licence, http://creativecommons.org/licenses/by/4.0/, as recorded in the arXiv licence metadata for this exact version and shown on the abstract page https://arxiv.org/abs/2609.09548v1. Full licence notice: \texttt{licenses/arxiv-2609.09548-CC-BY-4.0.txt}.\par\smallskip

\noindent\textit{Editorial scope.} Counted unit: the article's proposition eq:large\_l\_diag (source0.tex lines 2271--2295). The article's complete original proof of it is delivered with it (source0.tex lines 2297--2340, one \texttt{proof} environment of 44 lines). No mathematics is written, repaired or re-derived: the statement and the proof are the article's own lines, in the article's own notation, and no retained line is substituted or re-flowed.  Retained uncounted as the source-local context the proof needs: lines 1639--1651.  Nothing was omitted from any retained span, so the package carries no omission record.  No retained line contains a citation, so the wrapper carries no bibliography and no bibliography entry.  No retained line contains a figure environment, an image-inclusion command or drawing code, so no image is delivered and the package declares no asset dependency.  Editorial formatting only, in addition: an AMS \texttt{amsart} preamble that restates the article's own declarations and macros used by the retained lines, namely the environments \texttt{proposition} on separate counters, together with the standard packages those lines need.  The retained environments print in this document as Proposition 1 in order of appearance, whereas the article prints the same items with the numbering of its own full document.  The complete native source of the article as distributed in the arXiv e-print for this exact version is bundled unchanged as \texttt{sources/209827/source0.tex}.
\begin{align}
 G(x)
 &=2\sqrt{\frac2\pi}
 +\frac23\sqrt{\frac2\pi}\,x+O(x^2),
 \qquad x\downarrow0,
 \label{eq:Gsmall}
 \\
 G(x)
 &=\sqrt{2x}+\frac1{\sqrt{2x}}
 +O(\e^{-x}x^{-2}),
 \qquad x\to\infty.
 \label{eq:Glarge}
\end{align}

\begin{proposition}[Fixed radial indices, loss part]
For fixed $j$,
\begin{equation}
 \frac{\mathcal N^{(\ell)}_{jj}}{\kappa\pi}
 =
 \sqrt{2\beta}
 \left[
 1+
 \frac{3(2j+1)}{8\beta}
 +O(\beta^{-2})
 \right].
 \label{eq:large_l_diag}
\end{equation}
For fixed $d\ge1$,
\begin{equation}
 \frac{\mathcal N^{(\ell)}_{j,j+d}}{\kappa\pi}
 =
 \sqrt2
 \binom{1/2}{d}
 \sqrt{(j+1)_d}\,
 \beta^{(1-d)/2}
 \left[1+O(\beta^{-1})\right].
 \label{eq:large_l_offdiag}
\end{equation}
\end{proposition}

\begin{proof}
The diagonal matrix elements needed to relative order $\beta^{-1}$ are
\begin{equation}
 \langle j|T_\ell|j\rangle=2j,
\end{equation}
and
\begin{equation}
 \langle j|T_\ell^2|j\rangle
 =
 \beta(2j+1)+6j^2.
\end{equation}
For fixed radial indices this calculation is an \emph{entrywise}
asymptotic expansion, not a global operator-norm expansion of the unbounded
Jacobi operator.  Indeed, after compression to any fixed finite radial window,
$T_\ell/\beta=O(\beta^{-1/2})$ in finite-dimensional operator norm, so the
spectral-calculus expansion
\begin{equation}
 \sqrt{\beta I+T_\ell}
 =
 \sqrt\beta
 \left[
 I+\frac{T_\ell}{2\beta}
 -\frac{T_\ell^2}{8\beta^2}+\cdots
 \right]
\end{equation}
is legitimate to any prescribed order on that window.  For a fixed matrix
element, the coefficients through a prescribed order depend on only finitely
many Jacobi paths, so the same coefficients are those of the full operator.
Substitution of the two moments above, together with the leading
$J_\ell^{-1/2}$ contribution in \eqref{eq:Glarge}, yields
\eqref{eq:large_l_diag}.  For an off-diagonal
distance $d$, the first nonzero term occurs in $T_\ell^d$.  The unique minimal
raising path gives
\begin{equation}
 (T_\ell^d)_{j,j+d}
 =
 \prod_{r=0}^{d-1}
 \sqrt{(j+r+1)(\beta+j+r)}
 =
 \beta^{d/2}\sqrt{(j+1)_d}
 [1+O(\beta^{-1})].
\end{equation}
The binomial coefficient in \eqref{eq:large_l_offdiag} follows.
\end{proof}

\end{document}
